% Exercise 4.7
We have $K = 2$ classes, and if we let the stocks that issued dividends be
given by $k = 1$ and those that didn't by $k = 2$, we have $\pi_1 = 0.8$ and
$\pi_2 = 0.2$. Thus, from Bayes' theorem, we have that the probability that a
stock will issue a dividend, given given profit of $X = 4$, is given by the
following:
\begin{equation*}
    \begin{split}
        P(Y = 1 \vert X = 4) &= \frac{\pi_1 f_1(x)}{\sum_{l=1}^2 \pi_l f_l(x)} \\
        &= \frac{\pi_1 \frac{1}{\sqrt{2 \pi_1 \sigma^2}}e^{-{(4 - \bar{X})}^2 / 2\sigma^2}}
        {\sum_{l=1}^2 \pi_l \frac{1}{\sqrt{2 \pi_l \sigma^2}}
        e^{-{(4 - \bar{X}_l)}^2 / 2\sigma^2}} \\
        &= \frac{0.8 \frac{1}{\sqrt{1.6 \cdot 36^2}}e^{-{(4 - 10)}^2 / 2 \cdot 36^2}}
        {\sum_{l=1}^2 \pi_l \frac{1}{\sqrt{2 \pi_l \sigma^2}}
        e^{-{(4 - \bar{X}_l)}^2 / 2\sigma^2}} \\
        &= \frac{0.8 \frac{1}{\sqrt{1.6 \cdot 36^2}}e^{-{(4 - 10)}^2 / 2 \cdot 36^2}}
        {0.8 \frac{1}{\sqrt{1.6 \cdot 36^2}}e^{-{(4 - 10)}^2 / 2 \cdot 36^2}
        + 0.2 \frac{1}{\sqrt{0.4 \cdot 36^2}}e^{-{4^2 / 2 \cdot 36^2}}}.
    \end{split}
\end{equation*}
Thus, we have $P(Y = 1 \vert X = 4) \approx 0.6649498$, using the following
calculation:
\begin{lstlisting}
> pi_1 <- 0.8
> pi_2 <- 0.2
> xbar_1 <- 10
> xbar_2 <- 0
> sigma <- 36
> f <- function(x, pi, mu, sigma) {
+   (1 / sqrt(2 * pi * sigma^2)) * exp((-(x - mu)^2) / (2 * sigma^2))
+ }
> p <- (
+   pi_1 * f(4, pi_1, xbar_1, sigma) /
+     (pi_1 * f(4, pi_1, xbar_1, sigma) + pi_2 * f(4, pi_2, xbar_2, sigma))
+ )
> print(p)
[1] 0.6649498
\end{lstlisting}
